\chapter{Literature Review}

\section{Systemic Risk and Financial Contagion}

Systemic risk concerns the possibility that a disturbance initially
confined to one institution or market segment spreads through financial
linkages and feedback effects, eventually impairing the functioning of
the broader financial system. It differs from the stand-alone risk of an
individual bank because the scale of the final loss depends on how
shocks are transmitted and amplified after they occur. Capital buffers,
liquidity positions, and default probabilities remain important, but
they do not by themselves determine the system-wide outcome. Bilateral
claims, correlated asset holdings, and market responses also shape the
direction and intensity of risk propagation {[}7{]}.

Interbank claims are among the most direct channels of financial
contagion. Allen and Gale show that mutual claims can help banks share
regional liquidity shocks, although the same links can transmit losses
when liquidity becomes scarce. A network that supports risk sharing in
normal periods may therefore become a source of vulnerability under
stress. Gai and Kapadia develop this argument by examining contagion
across different network structures. Their results describe financial
systems as ``robust yet fragile'': many small shocks may be absorbed
without widespread failure, while a sufficiently large disturbance can
trigger a rapid cascade once critical counterparties become distressed.
Network connectivity should consequently be assessed together with the
size of exposures and the liquidity available to absorb losses.

Risk may also spread without a direct default relationship. Common asset
holdings can force institutions to recognize losses at the same time,
while fire sales depress market prices and weaken the balance sheets of
banks that did not initially experience a funding problem. Runs, credit
freezes, and declining confidence may reinforce these losses. Jackson
and Pernoud group these mechanisms into direct externalities and
feedback effects, emphasizing that systemic events usually emerge from
the interaction of several channels rather than from a single bilateral
exposure {[}6{]}. This broader perspective is relevant to interbank
markets because a payment shortfall can affect both the creditor's
balance sheet and its subsequent willingness or ability to provide
funding.

Network-based measures were developed partly because conventional
institution-level indicators do not capture these transmission effects.
DebtRank, proposed by Battiston et al., evaluates the potential distress
that an institution can transmit through its financial connections
{[}5{]}. The method illustrates why size alone is not a sufficient
measure of systemic importance. A small bank may still have a large
effect if it is linked to vulnerable counterparties or occupies a
position through which losses can reach several parts of the network.
Reviews of network models reach a similar conclusion: explicit claims,
common exposures, liquidity dependence, and other implicit connections
all influence the severity of systemic events {[}7, 8{]}.

These studies establish that systemic risk depends on both bank
conditions and the structure of financial connections. They also provide
methods for analyzing contagion after a set of bilateral exposures has
been specified. Less attention is given to the trading process through
which those exposures are created before clearing begins. In many
models, the network is treated as fixed or externally generated, even
though actual interbank relationships change with funding demand,
liquidity availability, interest-rate conditions, and counterparty
selection. This thesis addresses that limitation by generating interbank
loans endogenously under centralized matching and decentralized
request-for-quote mechanisms. The comparison examines whether different
trading processes produce different creditor--debtor networks and
whether those differences alter subsequent clearing losses, capital
stress, and systemic-risk dynamics.

\section{Interbank Network Models}

The interbank market can be represented as a directed and weighted
financial network. Banks form the nodes, while interbank loan contracts
define the edges between them. The direction of an edge identifies the
creditor and debtor, and its weight records the size of the
corresponding exposure. This representation captures both the
distribution of interbank claims and the channels through which payment
losses may later spread {[}17, 18{]}.

Empirical studies show that real interbank markets are neither random
nor uniformly connected. They are commonly characterized by sparse
links, heterogeneous exposures, and substantial differences in node
centrality {[}11, 18{]}. A limited number of institutions may act as
major funding providers or intermediaries, while most banks trade with a
limited set of counterparties. A bank's contribution to systemic risk
therefore depends not only on its balance-sheet condition, but also on
its position in the network and the size and direction of its exposures.

Many theoretical models treat the interbank network as an exogenous
input. Under this approach, researchers specify a complete, random,
sparse, or core--periphery structure and then examine how shocks
propagate through the imposed links. These studies show that
connectivity does not have a uniformly stabilizing effect. Additional
links may diversify small losses across several institutions, but they
can also transmit a sufficiently large shock to a wider part of the
system {[}4{]}. Network structure is therefore an important determinant
of contagion, although this approach does not explain how the structure
was formed.

In actual markets, lending relationships change as banks respond to
funding shortages, available liquidity, capital constraints,
interest-rate requirements, and counterparty risk. Previous transactions
may also affect the willingness of two banks to trade again,
particularly during periods of financial stress {[}10{]}. A fixed
network omits these decisions and assumes that the exposures through
which contagion occurs already exist. This limits its ability to
describe how risk concentration develops before payment difficulties
emerge.

Dynamic network models address this limitation by allowing bank
conditions and financial links to evolve together. New loans change
current liquidity positions and create future repayment obligations.
Outstanding contracts remain part of the network until they mature,
while repayment outcomes feed back into the resources available for
subsequent trading. The network observed in a particular period
therefore reflects both past lending decisions and current market
conditions. This intertemporal structure is relevant to systemic-risk
analysis because pressure may accumulate through repeated matching,
persistent exposures, scheduled repayments, and clearing losses.

The framework developed in this thesis follows this dynamic perspective.
The simulated banking system contains a central bank, commercial banks,
and shadow banks with different balance-sheet sizes, liquidity
positions, regulatory conditions, and risk preferences. Each active
institution may become a lender or borrower according to its current
funding position. Completed transactions are stored as dated loan
contracts containing the creditor, debtor, principal, interest rate, and
maturity information. The outstanding contracts are then aggregated into
the exposure and liability matrices required for network analysis and
clearing.

Interbank links are generated endogenously through two alternative
matching mechanisms. Centralized matching processes lending supply and
borrowing demand through a common market-wide allocation procedure. It
serves as a coordination benchmark in which available information and
funding opportunities can be considered jointly. Decentralized
request-for-quote matching instead limits borrowers to a set of
bilateral enquiries. Potential lenders decide whether and on what terms
to quote, after which the borrower selects among the available offers.
As a result, RFQ may leave part of the funding demand unmet even when
the system contains sufficient aggregate liquidity. It may also create
repeated relationships or concentrate exposures among institutions that
are contacted more frequently.

The difference between the two mechanisms concerns not only transaction
efficiency but also network formation. Centralized allocation can
distribute funds across a broader set of feasible counterparties,
whereas RFQ depends on local access, quotation decisions, and bilateral
selection. These differences can affect which banks obtain funding,
which lenders accumulate claims, and how repayment losses are
distributed when contracts mature. Holding the remaining balance-sheet
rules, contract accounting, and clearing procedure constant makes it
possible to relate later differences in systemic risk to the way the
creditor--debtor network was created.

Once scheduled obligations become due, the network shifts from loan
formation to settlement. A borrower that cannot meet its promised
payment reduces the amount received by its creditors, potentially
weakening their liquidity and capital positions. A loan contract enters
the risk process at several stages. Origination transfers liquidity
between the counterparties, the remaining principal contributes to
outstanding network exposure, and settlement may transmit losses if the
borrower cannot pay in full. Following the contract across these stages
links network formation directly to subsequent risk propagation.

Recent research has extended financial-network analysis beyond
single-layer interbank claims to bank--asset systems, bank--firm
relationships, and multilayer financial networks {[}19--24, 26{]}.
Despite differences in market setting, these studies reach a common
conclusion: systemic vulnerability reflects the interaction between
balance-sheet weakness and the structure of cross-institution exposures.
The present framework builds on this literature but differs in one
central respect. Rather than imposing a fixed network, it generates
interbank claims through centralized matching and decentralized RFQ,
preserves those claims as contracts, and follows them into later
settlement and systemic-risk measurement.

\section{Clearing Mechanisms}

Interbank loan contracts create mutually dependent payment obligations
between creditor and debtor banks. When one borrower cannot meet its
scheduled obligation in full, its creditors receive less than expected
and may in turn experience additional liquidity or capital pressure. A
clearing mechanism is therefore required to determine a payment vector
that is jointly feasible across the liability network rather than
evaluating each bank in isolation.

The Eisenberg--Noe framework provides a standard basis for this analysis
{[}2{]}. Its central contribution is to represent each bank's payment as
the smaller of its nominal obligation and the resources available from
external assets and incoming payments. Limited liability and
proportional repayment then determine how a shortfall is distributed
among creditors. Because incoming payments depend on the payments made
by other banks, the clearing outcome must be solved simultaneously for
the system as a whole.

Subsequent studies extend this framework by considering alternative
recovery assumptions, bankruptcy costs, intervention rules, and external
financial support {[}9{]}. These extensions show that contagion depends
not only on the liability network but also on the institutional rules
governing settlement and support. Nevertheless, the original framework
remains useful as a transparent benchmark for tracing how contractual
payment shortfalls propagate through interbank claims.

A dynamic model requires a further distinction between outstanding
exposure and payment currently due. An active loan may remain in the
network for several periods, while only the principal, installment, or
coupon scheduled for the current date should enter clearing. Treating
the entire outstanding balance as immediately payable would bring future
obligations forward and could cause the same contractual amount to be
processed more than once.

This thesis adopts the Eisenberg--Noe framework for due-payment clearing
after interbank claims have been generated through centralized matching
or decentralized RFQ. The same clearing rule is used under both
mechanisms, allowing differences in payment shortfalls and subsequent
risk dynamics to be related to the creditor--debtor structures created
during matching. The liability-matrix construction, fixed-point
solution, balance-sheet updates, and implementation of the dated
contract book are specified in Chapter~5.

\section{Liquidity and Solvency Regulation}

Capital and liquidity regulation provides an important link between
individual bank conditions and the stability of the wider financial
system. Capital determines the extent to which a bank can absorb losses,
whereas liquidity determines whether it can meet obligations when they
become due. These two dimensions are related but not interchangeable. A
bank may have sufficient capital yet encounter payment difficulties
because its assets cannot be converted into cash quickly enough.
Conversely, a bank with adequate short-term liquidity may remain
vulnerable if its capital buffer is too small to absorb losses.

The Basel III framework places particular emphasis on risk-sensitive
capital requirements. Rather than assessing capital relative to the
nominal size of a bank's balance sheet, the capital adequacy ratio
relates eligible capital to risk-weighted assets {[}14{]}. Assets with
different risk characteristics therefore make different contributions to
the capital requirement. This principle is especially relevant to
interbank network analysis because losses on interbank claims and
project investments can weaken a bank's capital position and reduce its
capacity to withstand subsequent shocks.

Liquidity regulation addresses a different source of vulnerability. The
regulatory liquidity coverage ratio compares high-quality liquid assets
with expected net cash outflows over a stressed period {[}12{]}. The net
stable funding ratio complements this short-horizon test by relating
available stable funding to required stable funding over a longer
horizon {[}13{]}. The present study operates at a daily frequency and
uses a simplified balance-sheet structure, so its liquidity measure
should be interpreted as a model-based proxy rather than a complete
implementation of either Basel III liquidity standard. The proxy
nevertheless captures the central economic relationship: liquidity
pressure rises when available liquid assets become insufficient relative
to expected funding outflows.

Capital and liquidity conditions can also influence banks before formal
default occurs. A bank facing a liquidity shortage may seek interbank
funding even when it remains solvent, while a bank with a weak capital
position may continue trading despite having limited capacity to absorb
further losses. Regulatory indicators therefore provide information
about emerging vulnerability that cannot be obtained from failure status
alone. This difference matters in a dynamic interbank network,
where payment shortfalls may weaken creditor banks before any
institution is classified as failed.

These conditions also shape market behavior. Banks with excess liquidity
may supply funds, whereas institutions facing short-term funding
pressure are more likely to borrow. Capital weakness can affect
resilience when interbank claims are not repaid in full, and settlement
losses may alter both capital and liquidity positions in later periods.
Regulation is therefore connected to network formation as well as risk
propagation: balance-sheet conditions influence trading decisions, while
the resulting lending relationships create exposures that feed back into
future balance sheets.

In the model, capital, liquidity, solvency, leverage, and activity status
are treated as evolving bank characteristics rather than static
regulatory observations. They are updated as market conditions,
investment outcomes, interbank transactions, and clearing results
change. The precise formulas, stress thresholds, system-level
indicators, and weighting scheme used in the simulation are specified in
Section~3.4.

\section{Matching Mechanisms in Interbank Markets}

Interbank networks are shaped not only by banks' balance-sheet
conditions but also by the process through which lenders and borrowers
meet. The matching mechanism determines which banks become
counterparties, how much funding is exchanged, and what contractual
terms are accepted. It therefore affects the number, direction, and size
of links in the resulting network. When these links are imposed in
advance, the analysis can examine contagion only after exposures already
exist. In the present framework, interbank relationships are instead
generated during the simulation, allowing network formation and
subsequent risk propagation to be studied within the same process.

The analysis compares centralized allocation with decentralized
request-for-quote matching. Both mechanisms use the same bank states,
regulatory constraints, contract-accounting rules, and clearing
procedure. Their main difference lies in how information is used and how
counterparties are selected. Centralized matching coordinates lending
supply and borrowing demand across the market, whereas RFQ limits each
borrower to a set of bilateral enquiries and quotations. Holding the
remaining modules fixed makes the matching rule the main source of
variation in funding outcomes and network structure.

\subsection{Centralized Matching Mechanism}

Centralized matching is used as the benchmark mechanism. At the
beginning of the matching stage, active banks are classified as lenders
or borrowers according to their liquidity positions and financing needs.
Each lender has a maximum amount of funds available for lending, while
each borrower has a specified funding demand.

The allocation procedure considers lenders and borrowers jointly. A
feasible loan plan is constructed subject to lender supply, borrower
demand, the upper limit on a single loan, network-degree restrictions,
interest-rate compatibility, and any applicable rollover constraints. A
borrower completes the transaction only when its demand can be satisfied
under these conditions. Otherwise, no new loan is created for that
borrower in the current period.

When a match is completed, the lender's liquid assets decrease and the
borrower's liquid assets increase by the transaction amount. The trade
is then stored as a dated contract containing the identities of the two
banks, the principal, the interest rate, and the maturity structure.
Because contracts remain active until their scheduled obligations are
settled, the interbank network is reconstructed from both newly issued
loans and exposures inherited from earlier periods.

Centralized matching represents a coordinated market environment.
Lending opportunities can be considered across the feasible market
rather than through a sequence of isolated bilateral enquiries. This
generally allows available liquidity to be allocated more efficiently
and may reduce unmet borrowing demand. The mechanism is nevertheless an
idealized benchmark, since real interbank markets do not always provide
complete information or a single process capable of coordinating all
transactions.

\subsection{Decentralized RFQ Matching Mechanism}

The decentralized mechanism follows a request-for-quote process. Banks
first form trading intentions from their current balance-sheet and
regulatory conditions. A lender's intended supply depends mainly on the
liquidity available above its target buffer, together with its
willingness to bear risk. A borrower's intended demand reflects its
funding gap, investment needs, and regulatory constraints. These
intentions convert individual bank conditions into potential lending
supply and borrowing demand.

A borrower does not submit its demand to a unified allocation process.
Instead, it contacts a limited number of candidate lenders and requests
quotations. Each lender decides whether to respond and what terms to
offer based on its remaining liquidity, minimum acceptable lending rate,
risk preference, and assessment of the borrower. The borrower then
selects among the quotations it has received.

A transaction is created only when an acceptable bilateral quote is
available. Completed trades are stored in the same contract book used
under centralized matching, so the subsequent accounting and clearing
procedures remain unchanged. The simulation also records funding
outcomes such as total lending supply, total borrowing demand,
transaction volume, funding satisfaction, unmet demand, the number of
completed trades, and average lending rates.

The RFQ mechanism changes access to funding even when aggregate market
liquidity is sufficient. A borrower may fail to obtain the required
amount because it contacts unsuitable lenders, receives quotations above
its acceptable rate, or is rejected because of counterparty risk.
Funding availability therefore depends on local access and bilateral
decisions rather than on system-wide supply alone.

This process can also produce a more uneven network. Banks that are
frequently contacted or have strong liquidity positions may accumulate a
larger number of claims, while other institutions remain weakly
connected. Repeated bilateral selection may create persistent trading
relationships or concentrate exposures among a small group of lenders.
These outcomes are not imposed in advance; they emerge from the enquiry
and quotation process.

\subsection{Comparison of Two Matching Mechanisms}

The central distinction between the two mechanisms concerns information
and counterparty access. Centralized matching evaluates feasible supply
and demand through a common allocation procedure. RFQ relies on limited
enquiries, lender-specific quotations, and borrower choice. The
difference is therefore not simply whether a loan is approved, but which
banks can reach each other and how lending capacity is distributed among
them.

Centralized matching considers a broader set of feasible counterparties,
but its all-or-nothing borrower rule may leave a borrower entirely
unfunded when complete coverage is unavailable. RFQ permits partial
funding but restricts access to a local candidate set. The relative
funding performance of the two mechanisms is therefore determined by the
simulation rather than assumed in advance.

RFQ introduces a different trade-off. Limited access and quotation
frictions may leave part of the borrowing demand unmet, causing
liquidity stress to appear earlier for some banks. However, the smaller
candidate set does not necessarily produce a less risky network. If the
same liquid banks are repeatedly selected, exposures may become
concentrated, and losses at settlement may be borne by a small group of
creditors.

The comparison is conducted under the same balance-sheet updates,
regulatory rules, contract treatment, and Eisenberg--Noe clearing
procedure. Any difference in later payment shortfalls or systemic-risk
indicators can therefore be related to the networks created by the
matching process rather than to different settlement assumptions. The
causal sequence is:


\begin{equation}
\begin{aligned}
\text{Bank states}&\rightarrow\text{lender and borrower roles}\rightarrow\text{matching}\rightarrow\text{loan contracts}\\
&\rightarrow\text{maturity payments}\rightarrow\text{clearing outcomes}\rightarrow\text{systemic risk}.
\end{aligned}
\end{equation}


Under this design, the matching mechanism affects systemic risk through
several channels. It changes the degree to which borrowing demand is
satisfied, determines which banks accumulate claims, alters the
concentration and maturity of exposures, and shapes the liability matrix
used during clearing. The experimental analysis therefore focuses on
whether centralized allocation and decentralized RFQ produce different
funding outcomes, network structures, and subsequent paths of liquidity
pressure, capital stress, and default contagion.

\section{Research Gap}

The literature establishes that contagion depends on both balance-sheet
capacity and the pattern of interbank claims, but these two strands are
often studied at different stages. Clearing models explain how losses
propagate through a given liability matrix, whereas empirical and
network studies describe sparsity, concentration, and centrality in
observed systems. Neither strand alone explains how a trading rule
creates the matrix that is later cleared.

A second gap concerns time. Static exposure matrices are useful for
stress tests, yet they do not distinguish a loan that remains
outstanding from a payment that is contractually due today. Once
maturities, installments, coupons, and rollover are introduced, this
distinction becomes necessary; otherwise the same principal can enter
clearing repeatedly or future obligations can be treated as immediately
payable.

The matching literature adds a further issue. Centralized allocation is
informative as a high-coordination benchmark, but over-the-counter
trading depends on limited enquiries and bilateral acceptance. Comparing
the two within one balance-sheet and clearing framework isolates a
practical question: does local trading reduce contagion by limiting
links, or increase realized failures because liquidity does not reach
some borrowers?

The present study addresses these gaps with a daily contract book. Bank
states generate supply and demand, the matching rule creates new
contracts, only scheduled payments enter the Eisenberg--Noe liability
matrix, and the clearing result feeds back into the next
day\textquotesingle s state. The contribution is therefore not a new
clearing theorem; it is an integrated setting in which network
formation, contract persistence, and clearing can be compared under
controlled simulation conditions.

The resulting framework combines bank-level capital and liquidity
conditions with endogenous network formation and due-payment clearing.
Centralized matching and decentralized RFQ are compared under the same
balance-sheet, contract-accounting, and clearing rules, while systemic
risk is evaluated through failures, capital stress, and capital
shortfalls. This design makes it possible to identify whether
differences in risk dynamics arise from the trading process that creates
interbank exposures.
